
\documentclass[11pt]{article}
\usepackage{amsmath,amssymb,amsthm,booktabs,geometry}
\geometry{margin=1in}
\newtheorem{theorem}{Theorem}
\newtheorem{definition}{Definition}
\title{Step 138: $\Omega$-Threshold Optimization for Blind-Projected Incidence}
\author{Six Birds RH Membrane Working Note}
\date{}
\begin{document}
\maketitle

\section{Purpose}
Steps 130--137 reduced the restricted source route to a metric-correct blind-sector estimate. The unrestricted GCD-log lower-frame import is blocked by squarefree Boolean near-null modes, but the actual residual coefficient class may avoid those modes after the zeta/Muentz incidence convolution. Step 138 chooses the threshold variables in a non-smuggled way.

\section{Definitions}
Let $P_y$ be a finite prime window and let
\[
(Z_y b)(T)=\sum_{S\subseteq T}b(S)
\]
be the zeta/Muentz incidence operator on squarefree seed coefficients. Let $\mathcal B_R$ be the declared blind Walsh family, for instance modes whose negative-prime set has cardinality at least $R$. Let $P_{\Omega\le K}$ be the upstream-declared seed cutoff.

\begin{definition}[Leakage and tail constants]
In the source-compatible central-line coefficient metric define
\[
L_{R,K}(y)=\|\Pi_{\mathcal B_R}Z_yP_{\Omega\le K}\|,
\qquad
T_{R,K}(y)=\|\Pi_{\mathcal B_R}Z_y(I-P_{\Omega\le K})\|.
\]
For the actual Burnol/Muentz residual seed window define
\[
\sigma_{K,N}=\|(I-P_{\Omega\le K})B_NG_{B,N}^{-1/2}\|.
\]
\end{definition}

\begin{theorem}[Threshold bound]
For every declared pair $(R,K)$,
\[
\|\Pi_{\mathcal B_R}Z_yB_NG_{B,N}^{-1/2}\|
\le L_{R,K}(y)+T_{R,K}(y)\sigma_{K,N}.
\]
Consequently
\[
\Xi_{\mathrm{GCD},N}\preceq \delta_{R,K,N}^2G_{B,N},
\qquad
\delta_{R,K,N}:=L_{R,K}+T_{R,K}\sigma_{K,N}.
\]
\end{theorem}

\begin{proof}
Decompose the residual seed into low and high $\Omega$ parts:
\[
B_N=P_{\Omega\le K}B_N+(I-P_{\Omega\le K})B_N.
\]
Apply $\Pi_{\mathcal B_R}Z_y$ and the triangle inequality. The first term is bounded by $L_{R,K}$ after normalization by $G_{B,N}^{-1/2}$; the second is bounded by $T_{R,K}\sigma_{K,N}$. Squaring gives the adequacy residual bound.
\end{proof}

\begin{theorem}[Restricted source-floor criterion]
Assume that the GCD-log source kernel has a lower frame on the nonblind sector,
\[
K_q\succeq \gamma_q\beta_R I
\quad\text{on }(I-\Pi_{\mathcal B_R})H,
\]
with $\gamma_q\to\infty$ and $\beta_R>0$. If
\[
\delta_{R,K,N}\le\delta_0<1,
\]
then on the residual coefficient class
\[
R_N^*K_qR_N\succeq \gamma_q\beta_R(1-\delta_0^2)G_{B,N}.
\]
Thus a positive blind-suppression floor is enough; exact suppression is sufficient but not necessary.
\end{theorem}

\section{Optimization rule}
For a declared admissible grid of thresholds, choose $(R,K,y)$ to maximize
\[
\mathfrak F(R,K,y)=\beta_R(y,q)\bigl(1-[L_{R,K}(y)+T_{R,K}(y)\sigma_{K,N}]^2\bigr)_+,
\]
subject to length, source, and fixed/exhaustive-tail records. This is a threshold-selection rule, not a proof of the seed-tail estimate.

\section{Nonclaim}
This step does not prove RH, the actual Burnol/Muentz seed tail estimate, or completed source coercivity. It supplies the finite threshold theorem that any such proof must pass.
\end{document}
